\begin{figure*}[t]
    \centering
    \includegraphics[width=\linewidth]{pdfs/sparse.pdf}\vspace{-3mm}
    \caption{{\bf ``Interacting" with an Image.} We translate a simple user input, mouse motions and drags, and expand it into a more complex motion prompt which helps to achieve the user's intention. The mouse trajectories are visualized as a hand when dragging, and as a black cursor otherwise. A grid of tracks centered on the cursor are created when the mouse is dragged, as shown in the top row. Frames from the generated video are shown in the bottom row. Prompting our model in this way, we can (a) move the head of a parrot or (c) a cow (b) play with hair or (d) ``interact" with an image of sand. We can also keep the background still by specifying static tracks, as in (b) or (d). \textbf{Note these samples are not generated in real-time and are not temporally causal. More examples can be found \supparxiv{in the \supp{}}{on our \href{https://motion-prompting.github.io/}{webpage}.}}
    }
\label{fig:sparse}
\vspace{2mm} %
\afterfig
\end{figure*}